\section{Measure full conversion calls}\label{rgb-measurement}
Both tested implementations use the same scalar pixel operations. The MSA version
adds an eight-pixel block path and a guard before the scalar tail. Measurements
include entry, block setup, conversion, tail and return. Caller preparation,
assembly, session creation, preview rendering and program exit are excluded.

\begin{table}[htbp]
  \centering
  \caption{Native instructions for the verified scalar and MSA implementations. Pixel count describes the same work for each pair of columns.}
  \tagpdfsetup{table/header-rows={1}}
  \begin{tabular}{rrrrr}
    \toprule
    Pixels & Scalar pack & MSA pack & Scalar expand & MSA expand \\
    \midrule
    0      & 5           & 8        & 5             & 8          \\
    1      & 23          & 26       & 30            & 33         \\
    7      & 131         & 134      & 180           & 183        \\
    8      & 149         & 54       & 205           & 55         \\
    9      & 167         & 72       & 230           & 80         \\
    16     & 293         & 76       & 405           & 86         \\
    31     & 563         & 224      & 780           & 292        \\
    32     & 581         & 120      & 805           & 148        \\
    \bottomrule
  \end{tabular}
\end{table}

Short calls pay for the MSA eligibility guard without executing a vector block.
A 31-pixel call needs three full blocks and seven scalar pixels; a 32-pixel call
needs four full blocks. The extra pixel can therefore reduce native work for these
implementations. State count and tail length when reporting such a comparison.

The vector path also repeats some byte accesses. For an eight-pixel pack, scalar
input reads cover 24 bytes once; the two vector loads produce 32 observed byte
reads while still remaining inside the same 24-byte slice. For expansion, scalar
output stores cover 24 bytes once and the two vector stores produce 32 observed
byte writes. Overlap makes the block fit; it does not make the extra accesses free.

In the pinned core, these vector transfers publish one memory notice per byte.
Observed bytes and notices therefore agree in this measurement, but neither counts
host memory-bus transactions. Native instruction reductions do not by themselves
establish elapsed-time speedups on the simulator's host. Exact counts are also not
free-writing grading requirements. Correct output, ownership and the public return
interface remain the comparison's starting conditions.
